\scanpage{441}
\mspage{8}
% archive sheet number: 440

\inserted{Soit $i\in\mathbb N$} On dit que $f$ est $i$-\uline{asphérique} si $\forall\,\unsure{t|s}\in\unsure{T|S}$, $Z=S/t$ est $i$-asphérique, \deleted{\illegible{1}} i.e. a $\pi_0(Z)\mathrel{-}\pi_i(Z)$ triviaux. On définit aussi $(-1)$-asphérique : les $S/t$ sont non vides. De \unsure{$=$|même} pour la notion de $i$-coasphérique.

\deleted{On doit avoir le\par \uline{Théorème ?} \inserted{Soit $f\colon S\longrightarrow T$ dans Cat, $i\in[\unsure{-0|0},\unsure{2|3}]=\{0,1,2,\deleted{\unsure{3}}\}$} a) Si $f$ est $(i-1)$-asphérique \inserted{et $(i-1)$-propre} ou \unsure{resp.} coasphérique \deleted{et de plus, lisse ou propre, (il suffit} $(i-1)$-lisse, \deleted{\unsure{ou $(i-1)$-propre})} alors le foncteur $f^*$ précédent est \deleted{\unsure{($i$}}\unsure{$i$}-fidèle. \par b) Supposons $f^*$ 2-fidèle, et soit $G$ un \deleted{\unsure{group}} fibré en groupoïdes sur $S$. Pour que $G$ soit $S$-équivalent à un \deleted{fibré en groupoïdes} $f^*(H)$, il \unsure{f.|faut} et s. que ce soit vrai \inserted{\unsure{colocalement|localement}} localement sur $T$}

En fait, pour $i\in[-1,+\infty[\ \subset\mathbb Z$, on définit un dérivateur $\unsure{\mathbb D|\mathrm D}_i$ des $i$-types d'homotopie relatifs (cf plus bas). D'où les notions d'objet $\mathbb D_i$-asphérique ($=i$-asphérique), de $\mathbb D_i$-équivalence, de foncteurs $\mathbb D_i$-asphériques et $\mathbb D_i$-coasphériques, de foncteurs $\mathbb D_i$-lisses et $\mathbb D_i$-propres, de $\mathbb D_i$-fibrations et de $\mathbb D_i$-\unsure{cofibrations|fibrations} \uline{triviales} \deleted{\unsure{dont}} (i.e. à fibres $\mathbb D_i$-asphériques). On doit avoir le

\uline{Théorème ?} Soit $f\colon S\longrightarrow T$ dans $(\mathrm{Cat})$, $i\in\{0,1,2\}$. \deleted{\unsure{a)}} \deleted{\unsure{Si $f$ est}} Supposons l'une des conditions suivantes \unsure{satisfaite|satisfaites} :

a) $f$ est $(i-1)$-propre, et $(i-1)$-asphérique (s'exprime par la condition que certaines catégories, telles les $S/t$, soient $(i-1)$-asphériques)

b) $f$ est $(i-1)$-lisse, et $(i-1)$-coasphérique.

Alors \textcircled{1} le foncteur $f^*$ d'image inverse

\[
f^*\colon\uline{\mathrm{Grp}}(T)\unsure{\mathrel{-}|\longrightarrow}\uline{\mathrm{Grp}}(S)
\]

est $i$-fidèle. \textcircled{2} Dans le cas $i=2$ ($f^*$ pl. fidèle) on
